\section{Interpretação das métricas}
O arrependimento é a diferença entre a maior utilidade do oráculo e a utilidade do candidato escolhido; zero indica uma escolha ótima segundo esse oráculo. A taxa de acerto é a fração de decisões que atingem esse máximo. A taxa de vitória ou empate nas ablações é a fração de decisões em que a atenção apresenta arrependimento menor ou igual ao comparador. Pearson $r$ mede associação linear, sem demonstrar causalidade ou independência quando próximo de zero. A semente fixa a sequência pseudoaleatória utilizada em cada coleta.
\section{Intervalos registrados e seus limites}
O relato por estimação é apoiado por \citet{cumming2014,wasserstein2016}. Nos scripts deste trabalho, a margem registrada é $1{,}96s_n/\sqrt{n}$, com $s_n$ calculado com divisor $n$. Com $n=200$ e margem $0{,}0105$, obtém-se $s_n\approx0{,}0758$. Essa identidade explica o cálculo preservado, sem validar a independência das decisões. A interpretação comparativa requer considerar o desenho, a dependência entre observações e o intervalo da diferença pareada. A não sobreposição de intervalos isolados não substitui essa análise.
Como cada piscina gera 40 decisões, a análise futura deve considerar o agrupamento por semente. Reamostragem \cite{efron1979} é uma possibilidade, desde que respeite a unidade de agrupamento; cinco piscinas ainda limitam a precisão. Não foi efetuado recálculo inferencial nesta revisão.
\section{Pareamento e calibração probabilística}
Na ablação com $q_w=0{,}6$, a diferença de arrependimento entre média igualitária e atenção é aproximadamente $+0{,}035$, com margem exploratória $0{,}011$. O sinal positivo descreve vantagem na amostra pareada registrada; não demonstra robustez fora desse desenho. A normalização de pesos não é calibração probabilística. Regras de escore próprias \cite{gneiting2007} oferecem fundamentos para outra análise, mas não foram calculadas curvas de calibração, erros de calibração ou escores probabilísticos nesta avaliação.
